% 11-protocole.tex — Résultats 1/4 : protocole d'évaluation (≈ 5 pages).

\chapter{Protocole d'évaluation}
\label{chap:protocole}

\section{Hypothèses et variables}

\begin{arediger}{1 page}
Reprendre les questions de recherche sous forme d'hypothèses testables, par
exemple : H1 les étudiants ayant utilisé les simulations progressent plus sur
les items de mouvement circulaire et de réaction normale ; H2 le mode jeu
augmente le temps d'utilisation volontaire ; H3 les formules substituées pas à
pas aident à relier calcul et phénomène. Variable indépendante, variables
dépendantes, variables contrôlées.
\end{arediger}

\section{Public et organisation}

\begin{arediger}{1,5 page}
Étudiants du cours de Physique~1 : effectifs, groupes, moment de l'année,
consentement et anonymisation. Plan expérimental (groupe expérimental et
groupe témoin, ou pré-test/post-test sur un même groupe) ; séances
d'utilisation (amphithéâtre, séances d'exercices, travail autonome).
\end{arediger}

\section{Instruments de mesure}

\begin{arediger}{2 pages}
Test conceptuel adapté (items inspirés du FCI, \cite{hestenes1992fci}, et
items construits sur le looping et le seau d'eau) ; questionnaire de
perception (échelle de Likert : utilité, clarté, motivation, charge
perçue) ; traces d'utilisation si elles sont disponibles (essais et
réussites du mode jeu) ; entretiens ou questions ouvertes. Validité et
fiabilité (alpha de Cronbach). Mettre les instruments complets en annexe.
\end{arediger}

\section{Méthodes d'analyse}

\begin{arediger}{½ page}
Gain normalisé de \textcite{hake1998}, tests statistiques adaptés à la taille
des groupes (test t apparié ou Wilcoxon), taille d'effet ; analyse
thématique des réponses ouvertes.
\end{arediger}
